% Figure 1 of the paper, for the project page. Build from this folder with ./build.sh,
% which also compiles a dark version by defining \darkmode before this file is read.
\documentclass[border=6pt]{standalone}
\usepackage[T1]{fontenc}
\renewcommand{\rmdefault}{ptm}
\usepackage{xcolor}
\usepackage{pgfplots}
\usepackage{pgfplotstable}
\pgfplotsset{compat=1.18}
\ifdefined\darkmode
  \definecolor{fg}{HTML}{E6E6E3}     % text and axes
  \definecolor{page}{HTML}{161719}   % the page behind the figure
  \definecolor{creal}{HTML}{8CB8E3}
  \definecolor{crand}{HTML}{F08C45}
  \definecolor{cnone}{HTML}{B5B5B0}
  \colorlet{bar}{fg!30!page}
  \colorlet{gridline}{fg!14!page}
  \colorlet{bracket}{fg!75!page}
\else
  \definecolor{fg}{HTML}{000000}
  \definecolor{page}{HTML}{FFFFFF}
  \definecolor{creal}{HTML}{1F4E79}
  \definecolor{crand}{HTML}{D55E00}
  \definecolor{cnone}{HTML}{4D4D4D}
  \colorlet{bar}{black!22}
  \colorlet{gridline}{black!10}
  \colorlet{bracket}{black!75}
\fi
\pgfplotsset{every axis/.append style={color=fg}}

\pgfplotstableread[comment chars={\#}]{figure-1.dat}\figdata
\newcommand{\cell}[2]{\pgfplotstablegetelem{#1}{#2}\of\figdata\pgfplotsretval}
% A number from the table, set in math mode so that its minus sign is a true minus.
\newcommand{\signed}[2]{\pgfplotstablegetelem{#1}{#2}\of\figdata
  $\pgfmathprintnumber[fixed, fixed zerofill, precision=1]{\pgfplotsretval}$}

\begin{document}
\color{fg}
\pgfplotstablegetelem{0}{rand}\of\figdata \edef\figrand{\pgfplotsretval}
\pgfplotstablegetelem{0}{real}\of\figdata \edef\figreal{\pgfplotsretval}
\pgfplotstablegetelem{0}{none}\of\figdata \edef\fignone{\pgfplotsretval}
\begin{tikzpicture}
\begin{axis}[
  width=8.4cm, height=6.6cm,
  xmin=0, xmax=60, ymin=-1.75, ymax=5.55, y dir=reverse,
  xlabel={Test Hits@20 (\%)},
  ytick={0,1,2,3,4,5},
  yticklabels={
    {GCN, 400 epochs}, {GCN, 4,000 epochs},
    {GAT, 400 epochs}, {GAT, 4,000 epochs},
    {GraphSAGE, 400 epochs}, {GraphSAGE, 4,000 epochs}},
  yticklabel style={font=\small},
  xticklabel style={font=\small},
  xlabel style={font=\small},
  axis x line*=bottom, axis y line*=left,
  y axis line style={draw=none}, ytick style={draw=none},
  xmajorgrids, grid style={gridline},
  clip=false,
  legend style={at={(0.5,1.03)}, anchor=south, draw=none, fill=none, legend columns=3,
    font=\small, column sep=6pt},
]
% The drop: a bar from the randomized-graph score to the real-graph score.
\addplot[forget plot, draw=bar, line width=3.4pt,
  quiver={u=\thisrow{real}-\thisrow{rand}, v=0}] table[x=rand, y expr=\coordindex]{\figdata};
\addplot[only marks, mark=o, mark size=2.7pt, line width=1.1pt, crand, mark options={fill=page}]
  table[x=rand, y expr=\coordindex]{\figdata};
\addlegendentry{randomized graph}
\addplot[only marks, mark=*, mark size=2.7pt, creal] table[x=real, y expr=\coordindex]{\figdata};
\addlegendentry{real graph}
\addplot[only marks, mark=square, mark size=2.9pt, line width=1.1pt, cnone]
  table[x=none, y expr=\coordindex]{\figdata};
\addlegendentry{no edges}
% The three quantities, marked above the first row.
\begin{scope}[bracket, line width=0.5pt, font=\scriptsize]
\draw (axis cs:\figrand,-1.12) -- (axis cs:\figrand,-1.22) -- node[above=-1pt]{penalty}
  (axis cs:\fignone,-1.22) -- (axis cs:\fignone,-1.12);
\draw (axis cs:\figrand,-0.45) -- (axis cs:\figrand,-0.55) -- node[above=-1pt]{drop}
  (axis cs:\figreal,-0.55) -- (axis cs:\figreal,-0.45);
\draw (axis cs:\figreal,-0.55) -- node[above=-1pt]{gain} (axis cs:\fignone,-0.55)
  -- (axis cs:\fignone,-0.45);
\end{scope}
% Drop and gain, printed beside each row.
\node[anchor=south, font=\small] at (axis cs:66,-0.5) {drop};
\node[anchor=south, font=\small] at (axis cs:86,-0.5) {gain, 95\% interval};
\pgfplotsinvokeforeach{0,...,5}{
  \node[font=\small] at (axis cs:66,#1) {\cell{#1}{drop}\%};
  \node[font=\small] at (axis cs:86,#1)
    {\signed{#1}{gain}\ [\signed{#1}{gain_low}, \signed{#1}{gain_high}]};
}
\end{axis}
\end{tikzpicture}
\end{document}
